\section{Toward Bethlehem}

\subsection{The House of Bread}\label{sec:bread}

Micah named the town, and the Fathers asked what its name meant. Bethlehem, in
the Hebrew, is the house of bread. Pope St Gregory the Great, preaching to the
people of Rome on the Gospel of Christmas night, made the name his sermon:
\latin{Bethlehem quippe domus panis interpretatur. Ipse namque est qui ait: Ego
sum panis vivus, qui de cælo descendi}. The older Roman Breviary read this homily
at Christmas Matins, and in its English the passage runs:

\begin{quote}
Now the name Bethlehem signifieth ``the House of Bread,'' and thus it is the
birth-place of Him Who hath said, ``I am the Living Bread, Which came down from
heaven.''\par
{\smallskip\noindent\small ---St Gregory the Great, Homily 8 on the Gospels\par}
\end{quote}

\noindent The place had borne the name long before, Gregory went on, because
the one who would feed the souls of his chosen was to appear there in the flesh.

Gregory saw a second sign in the road. The Child was born not in his parents'
house but on a journey, and, as the \emph{Catena Aurea} gives him, Gregory's reason was
this: ``And that He might shew that on account of the human form which He
took upon Him, He was born as in a strange country, not according to His power
but according to His nature.''

Then he came to the manger. Isaiah had said, in the chapter the Church reads on
the Tuesday of the second week of Advent, ``All flesh is grass.'' Gregory took
the prophet at his word: \latin{factus homo, fenum nostrum vertit in frumentum},
``made man, he turned our grass into wheat.'' And he went on: \latin{Unde et
natus in præsepio reclinatur, ut fideles omnes videlicet sancta animalia, carnis
suæ frumento reficeret}, ``So, when he is born, he is laid in a manger, that he
might refresh with the wheat of his flesh all the faithful, that is, the holy
animals.''

St Bede the Venerable took up the same thought, and carried it
from Bethlehem to every Christian soul:

\begin{quote}
But down to the very end of time, the Lord ceases not to be conceived at
Nazareth, to be born at Bethlehem, whenever any of His hearers taking of the
flour of the word makes himself a house of eternal bread. Daily in the Virgin's
womb, i.\,e. in the mind of believers, Christ is conceived by faith, born by
baptism.\par
{\smallskip\noindent\small ---St Bede, in the \emph{Catena Aurea} on Luke 2:6--7\par}
\end{quote}

\noindent St Cyril of Alexandria, as the \emph{Catena Aurea} gives him, said why
the Lord lay where the animals eat: so ``that we changing the life of beasts,
might be brought to the knowledge that befits man, partaking not of hay, but of
the heavenly bread, the lifegiving body.'' And St Gregory Nazianzen, on
Christmas Day, bade his people ``honour little Bethlehem, which hath led thee
back to Paradise; and worship the manger through which thou, being without
sense, wast fed by the Word.''

St Bernard, preaching on the vigil of Christmas to his monks, joined
the town's name to the name of its tribe: ``Bethlehem is the house of bread;
Juda signifies confession or praise.'' If a soul is fed with the Word, he
said, devoutly receives the Bread that came down from heaven, and lives by
faith, then ``you are become a Bethlehem fitted to receive our Lord.''

The crib in Christian homes owes much to St Francis of Assisi. Three
years before his death he kept Christmas at Greccio, in the hills of central
Italy, with a manger, hay, an ox and an ass, and St Bonaventure
told the story in his life of the saint:

\begin{quote}
The man of God, filled with tender love, stood before the manger, bathed in
tears, and overflowing with joy. Solemn Masses were celebrated over the manger,
Francis, the Levite of Christ, chanting the Holy Gospel. Then he preached unto
the folk standing round of the Birth of the King in poverty, calling Him, when
he wished to name Him, the Child of Bethlehem, by reason of his tender love for
Him.\par
{\smallskip\noindent\small ---St Bonaventure, \emph{Life of St Francis}\par}
\end{quote}

\noindent Pope Francis signed his apostolic letter on the crib, \emph{Admirabile
signum}, at Greccio, in its shrine of the Nativity, on 1 December 2019, the
First Sunday of Advent that year. Telling the story again after Thomas of
Celano, who wrote the first life of the saint, he noted that Greccio had no
carved figures: the friars and the country people who came were themselves the
scene, and one of them saw the Child lying in the manger (n.~2). The crib, he
wrote, shows how far the love of God would stoop, since the Maker of the world
became small for our sake; and he urged that a crib be set up at home, so that a
household may keep Bethlehem before its eyes, learn the poverty the Son of God
chose in becoming man, and serve him in the needy, as in St Matthew's account of
the judgment (n.~3). The first to come to the manger, he observed, were
shepherds, poor men keeping their flocks: the Incarnation was welcomed first by
the poor, and they keep a place of their own at the crib (nn.~5--6).
